\documentclass[profile=standard]{ipara-handbook}
\title{B04｜梯度、正交与 Lagrange：约束极值与子空间几何}
\author{}
\date{}
\begin{document}
\maketitle

\begin{quote}
本册回答：\textbf{为什么“约束极值”最后会出现梯度线性相关？Lagrange 乘子到底是在解一个公式，还是在描述“允许方向”与“禁止方向”的几何分解？}
\end{quote}

\section{本册定位：把约束先线性化，再用正交补读极值条件}
高数 Topic07 Own 梯度、约束极值和 Lagrange 方法；B01 Own 局部线性化；B00 Own 内积/正交；线代 Topic03 Own kernel、row space 与正交补。B04 只 Own 一条接口：
\[
\boxed{
\text{约束 }g(x)=c
\xrightarrow{\text{B01}}
Dg_x(h)=0
\xrightarrow{\text{kernel}}
T_xM=\ker Dg_x
\xrightarrow{\perp}
\text{法空间}
\xrightarrow{df|_{T_xM}=0}
\nabla f\text{ 落入法空间}
}
\]

这条链说明，Lagrange 乘子不是凭空出现的“新算法”，而是\textbf{约束线性化 + 正交补}的坐标表达。

\section{第一步：约束真正限制的是“允许小位移”}
设可行集
\[
M=\{x:g(x)=c\}.
\]
在一个正则点 $x$，若沿可行曲线 $\gamma(t)\subset M$ 运动且 $\gamma(0)=x$，则
\[
g(\gamma(t))\equiv c.
\]
对 $t$ 求导：
\[
Dg_x(\gamma'(0))=0.
\]
所以所有可行切向量都必须落在
\[
\boxed{\ker Dg_x}.
\]
在常见正则条件下，这个 kernel 正好就是约束曲面在 $x$ 处的切空间。

这一步非常重要：约束并不是“多了一条方程”，而是\textbf{删掉了一些可移动方向}。

\section{标量约束：为什么梯度就是法向}
若 $g:\mathbb R^n\to\mathbb R$ 可微，则
\[
Dg_x(h)=\nabla g(x)^T h.
\]
因此
\[
\ker Dg_x=\{h:\nabla g(x)^Th=0\}=\nabla g(x)^\perp.
\]
于是切空间与法方向形成正交分解：
\[
T_xM=\ker Dg_x,
\qquad
N_xM=\operatorname{span}\{\nabla g(x)\}.
\]

这就是高数“梯度垂直等值面”和线代“kernel 的正交补”真正接上的位置。

\section{极值条件：目标在所有允许方向上的一阶变化必须为零}
若 $x$ 是约束上的局部极值点，那么对任何可行切向 $h\in T_xM$，都必须有
\[
Df_x(h)=0.
\]
对标量目标函数
\[
Df_x(h)=\nabla f(x)^Th,
\]
所以
\[
\nabla f(x)\perp T_xM.
\]
而 $T_xM^\perp=N_xM$，因此
\[
\boxed{\nabla f(x)\in N_xM.}
\]
单个正则约束时法空间由 $\nabla g$ 张成，于是
\[
\boxed{\nabla f(x)=\lambda\nabla g(x).}
\]

所以“梯度平行”只是最终的坐标句；真正的母句是：
\[
\boxed{\text{目标的一阶变化在所有允许方向上都必须消失。}}
\]

\section{多约束：乘子其实是法空间坐标}
若
\[
g_1(x)=c_1,\dots,g_k(x)=c_k,
\]
把约束组合成
\[
g:\mathbb R^n\to\mathbb R^k.
\]
其 Jacobian $Dg_x$ 的 kernel 给切空间：
\[
T_xM=\ker Dg_x.
\]
若约束梯度线性无关，即 $Dg_x$ 满行秩，则
\[
T_xM^\perp=\operatorname{Im}(Dg_x)^T
=\operatorname{span}\{\nabla g_1,\dots,\nabla g_k\}.
\]
于是
\[
\boxed{
\nabla f(x)
=\lambda_1\nabla g_1(x)+\cdots+\lambda_k\nabla g_k(x).
}
\]

因此 Lagrange 乘子 $\lambda_i$ 可以理解为：\textbf{目标梯度在各约束法向基中的坐标。}

\section{为什么需要正则性：奇异约束会让“梯度法向”失去信息}
若单约束点满足
\[
\nabla g(x)=0,
\]
则线性化方程
\[
Dg_x(h)=0
\]
对所有 $h$ 都成立，无法正确刻画真实可行方向。

例如约束
\[
g(x,y)=x^2+y^2=0
\]
的可行集实际上只有原点，但在原点 $\nabla g=0$。若机械套
\[
\nabla f=\lambda\nabla g,
\]
就会要求 $\nabla f=0$，这显然不是“原点作为唯一可行点取得约束极值”的必要条件。

所以“约束梯度非零/多约束梯度独立”不是计算细节，而是保证\textbf{一阶线性化没有丢掉约束几何}的资格门。

\section{母例一：最近点为什么也是 Lagrange}
在直线
\[
x+y=1
\]
上找离原点最近点。令
\[
f(x,y)=x^2+y^2,
\qquad g(x,y)=x+y.
\]
切向方向满足
\[
h_1+h_2=0,
\]
即由 $(1,-1)$ 张成。约束极值要求目标梯度
\[
\nabla f=(2x,2y)
\]
对该切向正交，因此必须平行于法向 $(1,1)$：
\[
(2x,2y)=\lambda(1,1).
\]
结合 $x+y=1$ 得
\[
x=y=\frac12.
\]

这题同时可以由 B00 看成正交投影。两条路线共享“残差/梯度没有切向分量”这一结构，但一个研究点到子空间的距离，一个研究一般目标函数的约束驻点。

\section{母例二：为什么闭区域最值不能只解乘子方程}
Lagrange 只描述\textbf{光滑正则约束流形内部}的一阶候选。若题目是在闭区域上求全局最值，还必须把候选全集补齐：
\begin{itemize}
\item 内部驻点；
\item 每一条光滑边界上的正则约束候选；
\item 边界交点、端点、角点；
\item 约束奇异点；
\item 最后做函数值比较。
\end{itemize}

因此
\[
\boxed{\nabla f=\lambda\nabla g}
\]
从来不是“得到极值”的充分结论，它只是一类候选生成器。

\section{B04 串起哪些章节}
\begin{tabular}{p{0.23\linewidth}p{0.30\linewidth}p{0.39\linewidth}}
\hline
Owner / 下游 & 提供对象 & B04 的翻译责任\\
\hline
B01 / 高数 Topic07 & 约束的一阶线性化 $Dg_x$ & 把允许方向读成 $\ker Dg_x$\\
线代 Topic03 & kernel、row space、正交补 & 把切空间的正交补读成法空间\\
B00 / 线代 Topic01 & 内积、正交、投影 & 把“一阶变化为零”翻译成梯度无切向分量\\
高数 Topic07 & Lagrange 方程与全局最值 & 输出乘子候选的几何来源与资格边界\\
\hline
\end{tabular}

\section{更深一层：梯度本身依赖内积，微分才是第一对象}
$Df_x$ 本质上是一个线性泛函。之所以能够把它写成向量 $\nabla f$，是因为欧氏内积允许用
\[
Df_x(h)=\langle\nabla f,h\rangle
\]
把线性泛函表示成一个向量。

所以更稳定的抽象表达是
\[
Df_x|_{T_xM}=0.
\]
“梯度平行”是选定标准欧氏内积之后的向量化版本。这一背景能帮助阻断一个常见误区：换到非标准内积后，“梯度向量”的具体坐标会变化，但“微分在所有切向上为零”这条极值条件不变。

\section{反桥梁：最容易错误升级的结论}
\begin{tabular}{p{0.20\linewidth}p{0.04\linewidth}p{0.20\linewidth}p{0.48\linewidth}}
\hline
A & & B & 真正区别\\
\hline
$\nabla f\parallel\nabla g$ & $\neq$ & 已得到极值 & 只是正则约束下的一阶必要候选\\
$\nabla g=0$ & $\neq$ & 没有约束 & 可能是约束表示奇异，一阶线性化失真\\
切空间 & $\neq$ & 整个可行集 & 前者只是一点处的一阶线性近似\\
法空间 & $\neq$ & 一个固定法向量 & 多约束时法空间通常是多个约束梯度张成的子空间\\
\hline
\end{tabular}

\section{调用信号与退出边界}
题面出现“约束最值、最近点、切线/切平面、法向、多约束”时，先问：
\begin{enumerate}
\item 当前可行集在候选点是否正则？
\item 约束线性化 $Dg_x(h)=0$ 给出了哪些允许方向？
\item 目标的一阶变化是否在这些方向上全部为零？
\item 乘子候选之后是否还缺奇异点、端点、边界层级与全局比较？
\end{enumerate}

若题目进一步问约束下二阶充分条件，需要把 Hessian 限制到切空间上；这一层可作为 B03+B04 的 Extension，不在数学一主干机械展开。

\section{一页压缩}
B04 最后只保留一条母句：
\[
\boxed{
T_xM=\ker Dg_x,
\qquad
Df_x|_{T_xM}=0.
}
\]
在标准欧氏内积下，它等价于
\[
\boxed{\nabla f(x)\in T_xM^\perp.}
\]
单约束正则点再压缩成
\[
\nabla f=\lambda\nabla g.
\]

如果能从“允许方向”重新生成乘子公式，而不是从公式反背几何，B04 才真正掌握。

\end{document}
